\chapter{Settings and Preliminaries}\label{ch:preliminaries}

This chapter establishes the foundational framework for the free boundary min-max theory in complete Riemannian manifolds with boundary.

%%%%%%%%%%%%%%%%%%%%%%%%%%%%%%%%%%%%%%%%%%%%%%%%%%%%%%%%%%%%%%%%%%%%%%%%%%%%%%%
\section{Notation and Basic Conventions}\label{sec:notation}
%%%%%%%%%%%%%%%%%%%%%%%%%%%%%%%%%%%%%%%%%%%%%%%%%%%%%%%%%%%%%%%%%%%%%%%%%%%%%%%

Let $(M^{n+1}, g)$ be a smooth, complete Riemannian manifold with non-empty smooth boundary $\partial M$. We assume $2 \leq n \leq 6$ throughout. We write $\nu_{\partial M}$ for the inward unit normal to $\partial M$. See Appendix~\ref{sec:fermi-coordinates} for details on Fermi coordinates near $\partial M$.

We denote by $\mathcal{H}^k$ the $k$-dimensional Hausdorff measure and by $\mathbf{M}(T)$ the mass of a current $T$. For vector fields, we write $\mathfrak{X}(M)$ for smooth compactly supported vector fields on $M$, and 
\[
\mathfrak{X}_{\mathrm{tan}}(M) := \{X \in \mathfrak{X}(M) : X(p) \in T_p(\partial M) \text{ for all } p \in \partial M\}
\]
for those tangent to $\partial M$.

For a relatively open subset $U \subset M$, we write $\mathrm{int}_M(U)$ for the relative interior of $U$ in $M$, and $\partial_{\mathrm{rel}} U$ for the relative boundary. We use $A \Subset B$ to denote that $\overline{A}$ is compact and contained in $B$.

%%%%%%%%%%%%%%%%%%%%%%%%%%%%%%%%%%%%%%%%%%%%%%%%%%%%%%%%%%%%%%%%%%%%%%%%%%%%%%%
\section{Currents and Relative Cycles}\label{sec:relative-cycles}
%%%%%%%%%%%%%%%%%%%%%%%%%%%%%%%%%%%%%%%%%%%%%%%%%%%%%%%%%%%%%%%%%%%%%%%%%%%%%%%

We work with currents and relative cycles following Li--Zhou~\cite{LZ17}. We write $\mathcal{R}_k(M)$ for the space of integer rectifiable $k$-currents supported in $M$, with mass norm $\mathbf{M}$ and flat topology. For a bounded open set $\Omega \subset M$ of finite perimeter (i.e.\ a Caccioppoli set), we write $Q_\Omega := \llbracket \Omega \rrbracket \in \mathcal{R}_{n+1}(M)$ for the associated $(n+1)$-current, so that $\mathbf{M}(\partial Q_\Omega) < \infty$.

\begin{definition}[Relative cycles]\label{def:relative-cycle}
The space of \emph{relative cycles} is
\[
\mathcal{Z}_n(M, \partial M) := \{T \in \mathcal{R}_n(M) : \mathrm{spt}(T) \subset M, \ \mathrm{spt}(\partial T) \subset \partial M\},
\]
equipped with the flat topology. Each $\tau \in \mathcal{Z}_n(M, \partial M)$ satisfies $\partial \tau \subset \partial M$ automatically.
\end{definition}

\begin{remark}\label{rem:filling-cycle}
For any bounded open set $\Omega \subset M$ of finite perimeter, the current $\partial \llbracket \Omega \rrbracket$ gives a relative cycle in $\mathcal{Z}_n(M, \partial M)$. We refer to $\llbracket \Omega \rrbracket$ as the \emph{filling} of this cycle. For a relative cycle $\tau \in \mathcal{Z}_n(M, \partial M)$, the mass restricted to an open set $V$ is $\|\tau\|(V) := \mathbf{M}(\tau \llcorner V)$.
\end{remark}

%%%%%%%%%%%%%%%%%%%%%%%%%%%%%%%%%%%%%%%%%%%%%%%%%%%%%%%%%%%%%%%%%%%%%%%%%%%%%%%
\section{Boundary Decomposition of Currents}\label{sec:boundary-decomposition}
%%%%%%%%%%%%%%%%%%%%%%%%%%%%%%%%%%%%%%%%%%%%%%%%%%%%%%%%%%%%%%%%%%%%%%%%%%%%%%%

For a bounded open set $U \subset M$ with smooth boundary, the associated $(n+1)$-current $\llbracket U \rrbracket \in \mathcal{R}_{n+1}(M)$ has boundary
\[
\partial \llbracket U \rrbracket = \partial_I U + \partial_F U,
\]
where:
\begin{itemize}
    \item $\partial_I U := \partial \llbracket U \rrbracket \llcorner \mathrm{int}(M) \in \mathcal{Z}_n(M, \partial M)$ is the \emph{interior boundary} (a relative cycle);
    \item $\partial_F U := \partial \llbracket U \rrbracket \llcorner \partial M$ is the \emph{free boundary part} (supported on $\partial M$).
\end{itemize}

%%%%%%%%%%%%%%%%%%%%%%%%%%%%%%%%%%%%%%%%%%%%%%%%%%%%%%%%%%%%%%%%%%%%%%%%%%%%%%%
\section{Sweepouts of Open Sets}\label{sec:sweepouts}
%%%%%%%%%%%%%%%%%%%%%%%%%%%%%%%%%%%%%%%%%%%%%%%%%%%%%%%%%%%%%%%%%%%%%%%%%%%%%%%

\begin{definition}[Sweepout]\label{def:sweepout}
Let $U \subset M$ be a bounded open set of finite perimeter. A \emph{sweepout of $U$} is a family $\{\tau_t\}_{t \in [0,1]}$ of relative cycles $\tau_t \in \mathcal{Z}_n(M, \partial M)$ together with a family of fillings $\{Q_t\}_{t \in [0,1]}$, $Q_t \in \mathcal{R}_{n+1}(M)$ with $\partial Q_t = \tau_t + S_t$ where $S_t$ is supported on $\partial M$, satisfying:
\begin{enumerate}[label=\textup{(sw\arabic*)}]
    \item\label{sw:finite} $\mathbf{M}(Q_t) < \infty$ for all $t$;
    \item\label{sw:endpoints} $Q_0 \llcorner U = 0$ and $Q_1 \llcorner U = \llbracket U \rrbracket$;
    \item\label{sw:continuity} The map $t \mapsto Q_t$ is continuous in the flat topology;
    \item\label{sw:no-concentration} (No mass concentration) For every $x \in M$,
    \[
    \limsup_{r \to 0} \sup_{t \in [0,1]} \|\tau_t\|(B_r(x)) = 0.
    \]
\end{enumerate}
\end{definition}

\begin{remark}\label{rem:free-boundary-automatic}
Since each $\tau_t \in \mathcal{Z}_n(M, \partial M)$, the free boundary condition $\partial \tau_t \subset \partial M$ is automatic.
\end{remark}

\begin{definition}[Height and width]\label{def:height-width}
The \emph{height} of a sweepout $\{\tau_t\}$ is $\mathcal{F}(\{\tau_t\}) := \sup_{t \in [0,1]} \mathbf{M}(\tau_t)$.
The \emph{width} of a collection $\mathcal{S}$ of sweepouts is $W(\mathcal{S}) := \inf_{\{\tau_t\} \in \mathcal{S}} \mathcal{F}(\{\tau_t\})$.
\end{definition}

%%%%%%%%%%%%%%%%%%%%%%%%%%%%%%%%%%%%%%%%%%%%%%%%%%%%%%%%%%%%%%%%%%%%%%%%%%%%%%%
\section{Nested and Good Sweepouts}\label{sec:nested-good}
%%%%%%%%%%%%%%%%%%%%%%%%%%%%%%%%%%%%%%%%%%%%%%%%%%%%%%%%%%%%%%%%%%%%%%%%%%%%%%%

\begin{definition}[Nested sweepout]\label{def:nested-sweepout}
A sweepout $\{\tau_t, Q_t\}$ of $U$ is \emph{nested} if:
\begin{enumerate}[label=\textup{(n\arabic*)}]
    \item\label{n:monotone} $Q_s \leq Q_t$ for all $s \leq t$ (i.e.\ $Q_t - Q_s$ is a non-negative current);
    \item\label{n:morse} There exists a proper Morse function $f: M \to \mathbb{R}$ with $f$ transverse to $\partial M$, such that $Q_t = \llbracket f^{-1}((-\infty, c_t)) \rrbracket$ for some continuous increasing function $t \mapsto c_t$.
\end{enumerate}
We denote the collection of nested sweepouts of $U$ by $\mathcal{S}_n(U)$, and write $W_n(U) := W(\mathcal{S}_n(U))$.
\end{definition}

\begin{definition}[Good sweepout]\label{def:good-sweepout}
A sweepout $\{\tau_t, Q_t\}$ of $U$ is a \emph{good sweepout} if it additionally satisfies:
\begin{enumerate}[label=\textup{(g\arabic*)}]
    \item\label{g:endpoint-bound} $\mathbf{M}(\tau_0) \leq 5\,\mathbf{M}(\partial_I U)$ and $\mathbf{M}(\tau_1) \leq 5\,\mathbf{M}(\partial_I U)$.
\end{enumerate}
We denote the collection of good sweepouts of $U$ by $\mathcal{S}_g(U)$, and write $W_g(U) := W(\mathcal{S}_g(U))$.
\end{definition}

\begin{definition}[Relative sweepout]\label{def:relative-sweepout}
Given a nested sweepout $\{\tau_t, Q_t\}$ of $U$, the \emph{relative sweepout} is the family
\[
\sigma_t := \tau_t \llcorner \overline{U} \in \mathcal{Z}_n(\overline{U}, \partial M \cap \overline{U}).
\]
We denote the collection of relative sweepouts of $U$ by $\mathcal{S}_\partial(U)$, and write $W_\partial(U) := W(\mathcal{S}_\partial(U))$.
\end{definition}

\begin{proposition}[Width inequalities]\label{prop:width-inequalities}
Let $U \subset M$ be a bounded open set with smooth boundary. Then:
\begin{enumerate}[label=\textup{(\roman*)}]
    \item $W_\partial(U) \leq W_n(U) \leq W_\partial(U) + \mathbf{M}(\partial_I U)$;
    \item $W_n(U) = W(U)$, where $W(U)$ is the width over all sweepouts;
    \item If $U$ is a good set (Definition~\ref{def:good-set}), then $W_g(U) = W(U)$.
\end{enumerate}
\end{proposition}

%%%%%%%%%%%%%%%%%%%%%%%%%%%%%%%%%%%%%%%%%%%%%%%%%%%%%%%%%%%%%%%%%%%%%%%%%%%%%%%
\section{Good Sets}\label{sec:good-sets}
%%%%%%%%%%%%%%%%%%%%%%%%%%%%%%%%%%%%%%%%%%%%%%%%%%%%%%%%%%%%%%%%%%%%%%%%%%%%%%%

\begin{definition}[Good set]\label{def:good-set}
A bounded open subset $U \subset M$ of finite perimeter is a \emph{good set} if
\[
\mathbf{M}(\partial_I U) \leq \frac{1}{10}\, W_\partial(U),
\]
where $\partial_I U = \partial \llbracket U \rrbracket \llcorner \mathrm{int}(M)$ is the interior boundary of $U$.
\end{definition}

\begin{remark}[The bottleneck principle]\label{rem:bottleneck}
The good set condition ensures that mass cannot ``escape to infinity'' through $\partial_I U$. Any sweepout of $U$ must move the filling from $Q_0 \llcorner U = 0$ to $Q_1 \llcorner U = \llbracket U \rrbracket$. This transfer must pass through the ``bottleneck'' $\partial_I U$, whose mass satisfies $\mathbf{M}(\partial_I U) \ll W_\partial(U)$. Therefore at least one slice $\tau_t$ must carry substantial mass inside $\overline{U}$.

This principle is made rigorous in Chapter~\ref{ch:no-escape-to-infinity} (Proposition~\ref{prop:non-escape}).
\end{remark}

\begin{lemma}[Existence of good sets]\label{lem:good-set-existence}
Let $(M^{n+1}, g)$ be a complete Riemannian manifold with non-empty boundary and sublinear volume growth, i.e.\
\[
\liminf_{R \to \infty} \frac{\mathcal{H}^{n+1}(B_R(x) \cap M)}{R} = 0 \qquad \text{for some } x \in M.
\]
Then there exists a good set $U \subset M$ with $0 < W_g(U) < \infty$.
\end{lemma}

\begin{proof}
Fix a point $x \in M$ and a radius $r > 0$ such that the local isoperimetric constant
\[
C_I := \inf\bigl\{\mathbf{M}(\partial_I E) : E \subset B_r(x),\ \mathcal{H}^{n+1}(E) = \tfrac{1}{2}\,\mathcal{H}^{n+1}(B_r(x))\bigr\} > 0.
\]
Let $d_x: M \to [0, \infty)$ denote the distance function from $x$. Applying the slicing theorem (Theorem~\ref{thm:slicing}) to the current $T = \llbracket M \rrbracket$ and the Lipschitz function $f = d_x$ with $|\nabla d_x| \leq 1$, we obtain
\[
\int_0^\infty \mathbf{M}(\langle \llbracket M \rrbracket, d_x, R \rangle)\, dR \leq \mathcal{H}^{n+1}(M \cap B_S(x))
\]
for any $S > 0$. The sublinear growth assumption therefore allows us to pick a regular value $R > r$ such that
\[
\mathbf{M}(\langle \llbracket M \rrbracket, d_x, R \rangle) \leq \frac{C_I}{100}.
\]
Set $U := B_R(x) \cap M$. Since $\partial_I U$ is supported on $\partial B_R(x) \cap \mathrm{int}(M)$, we have $\mathbf{M}(\partial_I U) \leq \mathbf{M}(\langle \llbracket M \rrbracket, d_x, R \rangle) \leq C_I/100$. Any sweepout of $U$ must bisect volume in $B_r(x)$ at some time, so $W_\partial(U) \geq C_I$. Since $\mathbf{M}(\partial_I U) \leq \frac{1}{10} W_\partial(U)$, the set $U$ is good. Finiteness $W_g(U) < \infty$ follows from the existence of a sweepout by slices $\langle \llbracket M \rrbracket, d_x, R' \rangle$ for $R' \in [0, R]$.
\end{proof}


%%%%%%%%%%%%%%%%%%%%%%%%%%%%%%%%%%%%%%%%%%%%%%%%%%%%%%%%%%%%%%%%%%%%%%%%%%%%%%%
\section{Slicing of Currents}\label{sec:slicing}
%%%%%%%%%%%%%%%%%%%%%%%%%%%%%%%%%%%%%%%%%%%%%%%%%%%%%%%%%%%%%%%%%%%%%%%%%%%%%%%

We recall the slicing theory for integer rectifiable currents; see Simon~\cite[Sections~28--29]{Sim84}.

\begin{theorem}[Slicing~{\cite[28.4, 28.5]{Sim84}}]\label{thm:slicing}
Let $T \in \mathcal{R}_{n+1}(M)$ and let $f: M \to \mathbb{R}$ be a Lipschitz function. Then for almost every $t \in \mathbb{R}$, the \emph{slice}
\[
\langle T, f, t \rangle \in \mathcal{R}_n(M)
\]
is a well-defined integer rectifiable $n$-current satisfying:
\begin{enumerate}[label=\textup{(\roman*)}]
    \item (Boundary formula) $\partial \langle T, f, t \rangle = \langle \partial T, f, t \rangle$ wherever $\partial T$ has locally finite mass;
    \item (Coarea inequality) $\int_{\mathbb{R}} \mathbf{M}(\langle T, f, t \rangle)\, dt \leq \int |\nabla f|\, d\|T\|$.
\end{enumerate}
\end{theorem}

\begin{corollary}[Coarea formula for fillings]\label{cor:coarea-filling}
Let $Q = \llbracket \Omega \rrbracket$ for a bounded open set $\Omega \subset M$ of finite perimeter, and let $f: M \to \mathbb{R}$ be Lipschitz. Then for almost every $t$, the slice $\langle Q, f, t \rangle \in \mathcal{Z}_n(M, \partial M)$ is a relative cycle, and
\[
\int_{\mathbb{R}} \mathbf{M}(\langle Q, f, t \rangle)\, dt \leq \int_\Omega |\nabla f|\, d\mathcal{H}^{n+1}.
\]
\end{corollary}